\documentclass{article}
\usepackage[nohead]{geometry}
\usepackage{amsmath}
\usepackage{verbatim,fancyvrb}
\usepackage{color,gretl}
\usepackage{graphicx}
\usepackage[authoryear]{natbib}
\usepackage{smallhds,url}
\usepackage{hyperref}

\setlength{\parindent}{0pt}
\setlength{\parskip}{1ex}

\begin{document}

\begin{center}
  {\Large \textbf{holt\_winters 0.1}} \\[6pt]
  Allin, 2026-03-09
\end{center}

\section{Introduction}

This is a draft of a package to implement the Holt--Winters
forecasting method, namely triple exponential smoothing with level,
trend and seasonal components.

The package requires time-series data with a seasonal dimension. Let
$m$ denote the periodicity of the data: 4 for quarterly data; 12 for
monthly; 5, 6 or 7 for daily; 24 for hourly.

The notation used here is based on Rob Hyndman's exposition at
\url{https://otexts.com/fpp3/holt-winters.html}. There are two
variants of Holt--Winters, differing with respect to their treatment
of the seasonal component, which can be either multiplicative or
additive. The multiplicative variant is the default; it is represented
by the following recursion, in which $\ell$ denotes the level, $b$ the
trend or slope, and $s$ the seasonal
\begin{align*}
  \ell_t &= \alpha (y_t / s_{t-m}) + (1-\alpha) (\ell_{t-1} + b_{t-1})\\
  b_t &= \beta (\ell_t - \ell_{t-1}) + (1-\beta) b_{t-1}\\
  s_t &= \gamma y_t / \ell_t + (1-\gamma) s_{t-m}\\
  \hat{y}_{t+h|t} &= (\ell_t + hb_t)\, s_{t+h-m(k+1)}
\end{align*}
where $h$ denotes the number of steps ahead for the forecast and $k$
is the integer part of $(h-1)/m$ ($=0$ when when $h=1$).

\section{The public function}

The \dtk{holt_winters()} function has this signature:

\begin{code}
  function bundle holt_winters (series y,
                                matrix theta[auto],
                                matrix init[auto],
                                bool additive[0],
                                int os_tail,
                                int verbose[0:3:1)
\end{code}

Description of the arguments:

\begin{itemize}
  \item \texttt{y} (required): the series for forecasting.
  \item \texttt{theta} (optional): if supplied, this must be a
    3-vector holding the smoothing parameters $\theta$ = ($\alpha$,
    $\beta$, $\gamma$), with $0 \le \theta_i \le 1,\, \forall i$.  If
    \texttt{parms} is not given these parameters are determined
    automatically.
  \item \texttt{init} (optional): if supplied, this must be a vector
    of length $2+m$ holding initial values. Let $\tau$ indicate the
    first observation for which a forecast is wanted; then the
    \texttt{init} values should be $\ell_{\tau-1}$, $b_{\tau-1}$, and
    $s_{\tau-i}, i = m, m-1, \dots,1$. If \texttt{init} is not given
    the initializer is set automatically.
  \item \texttt{additive} (optional): is a boolean flag to specify
    that the seasonal should be treated as additive rather than
    multiplicative.
  \item \dtk{os_tail} (required): this specifies the number of
    observations to reserve for out-of-sample forecasting. 
  \item \texttt{verbose}: controls the verbosity of output, with four
    levels. \texttt{verbose = 0} makes the function operate silently.
\end{itemize}

\section{Automatic usage}
\label{sec:auto}

\section{Rule-of-thumb initialization}
\label{sec:init}

\section{The additive method}
\label{sec:additive}

\end{document}
